\documentclass[11pt,a4paper]{article}
\usepackage[utf8]{inputenc}
\usepackage[T1]{fontenc}
\usepackage[spanish,es-nodecimaldot]{babel}
\usepackage{amsmath,amssymb}
\usepackage{booktabs}
\usepackage[margin=2.4cm]{geometry}
\usepackage{hyperref}
\usepackage{xcolor}
\hypersetup{colorlinks=true,linkcolor=blue!50!black,urlcolor=blue!50!black}
\usepackage{fancyvrb}
\usepackage{titlesec}
\titleformat{\section}{\large\bfseries}{\thesection.}{0.6em}{}
\setlength{\parskip}{0.5em}
\setlength{\parindent}{0pt}

\title{\bfseries La consulta no puede ver la pregunta\\[0.3em]
\large El alcance de una convolución corta decide qué parte de una pregunta\\
condiciona la búsqueda en la memoria de un modelo}
\author{Maximiliano Speranza\\
\small Investigador independiente, Buenos Aires, Argentina\\
\small \href{https://orcid.org/0009-0005-0413-8554}{ORCID 0009-0005-0413-8554}}
\date{Septiembre de 2026}

\begin{document}
\maketitle

\begin{center}
\fbox{\begin{minipage}{0.92\textwidth}
\small
\textbf{Resumen en una línea.} En los modelos de atención lineal y de espacio de estados hay una parte
de cada pregunta que la búsqueda interna \textbf{literalmente no puede ver}, y lo que queda afuera no
se degrada de a poco: \textbf{desaparece exacto}. El modelo entonces contesta con confianza usando
sólo la parte que sí alcanza a ver.
\end{minipage}}
\end{center}

\section{Qué es}

Un modelo con memoria tiene que hacer dos cosas que suelen medirse juntas y son distintas. Tiene que
\emph{encontrar} lo que busca, y tiene que \emph{saber} cuándo lo que busca no está. La segunda es la
que separa una memoria de un fabulador.

Lo que encontramos es que la segunda puede fallar por una causa puramente mecánica, aguas arriba de
todo lo que el modelo aprende. La consulta con la que se busca en la memoria no se arma con la
pregunta entera. Se arma con una \textbf{convolución causal corta} sobre los estados ocultos, y esa
convolución tiene un alcance fijo de pocos tokens.

Si una parte de la pregunta cae afuera de ese alcance, no influye poco en la búsqueda.
\textbf{No influye nada.}

\section{Cómo funciona}

El tamaño de núcleo por omisión de esa convolución es \textbf{4}. Es decir que la consulta de una capa
es función del token actual y de \textbf{tres anteriores}, y de nada más.

Ese número no suele medirse: se hereda. Mamba-2 usa 4 por defecto. La implementación de referencia de
Qwen3-Next expone \texttt{linear\_conv\_kernel\_dim}, documentado como «kernel size of the convolution
used in linear attention layers», también con valor 4. Un alcance de tres tokens no es una
configuración exótica, es el ajuste de fábrica en modelos desplegados.

\subsection*{La falla concreta que produce}

Pensemos una pregunta de la forma \emph{«cuál es la \textbf{⟨relación⟩} de \textbf{⟨entidad⟩}?»}.
Leída desde el final, la \textbf{entidad} queda a 1 token y la \textbf{relación} a 3. Con alcance 2,
la relación cae \textbf{exactamente un token afuera}, de forma determinista, en el 100\,\% de las
consultas.

El modelo busca entonces \textbf{usando sólo la entidad}. Y ahí aparece la falla que más se parece a
una alucinación real: si se le pregunta por una relación que \textbf{nunca se enunció}, sobre una
entidad que \textbf{sí} existe, no se abstiene. Recupera el hecho vecino, el que sí tiene guardado de
esa entidad, y lo entrega con confianza.

No estaba ignorando la relación. \textbf{No la podía ver.}

\section{Cómo se descubrió}

La intervención es tan simple que se puede repetir en cualquier arquitectura, y no necesita entrenar
nada:

\begin{center}
\fbox{\begin{minipage}{0.85\textwidth}
\centering\textbf{Cambiá un token de la pregunta y mirá si la búsqueda se mueve.}
\end{minipage}}
\end{center}

Se toma una posición de lectura, se cambia \textbf{un solo token} a distancia $d$ hacia atrás, y se
mide cuánto se movió la salida de la convolución en esa posición.

\subsection*{Medición 1 · el alcance real no es el nominal}

Antes de medir el efecto hay que medir la ventana, porque un peso puede estar en cero y entonces el
alcance efectivo es menor que el nominal. En \texttt{mamba-130m-hf}:

\begin{center}
\begin{tabular}{crrrr}
\toprule
capa & \multicolumn{4}{c}{$\max|\text{peso}|$ de cada tap, del más viejo al actual} \\
\midrule
0  & \textbf{0.000000} & 0.933774 & 1.396596 & 3.513000 \\
1  & \textbf{0.000000} & 0.599272 & 0.928984 & 1.315802 \\
12 & \textbf{0.000000} & 0.309363 & 0.901561 & 0.981588 \\
23 & \textbf{0.000000} & 0.128891 & 0.611719 & 1.090700 \\
\bottomrule
\end{tabular}
\end{center}

\textbf{El tap más viejo es cero exacto en las 24 capas.} No pequeño: cero, en las 1.536 entradas de
cada capa. El alcance efectivo de este checkpoint es \textbf{2 tokens}, no 3.

Informamos la medición y no una causa. Un peso exactamente nulo en $24 \times 1536$ entradas no sale
de un gradiente, así que un artefacto de conversión es plausible, pero no lo verificamos y no lo
afirmamos. La consecuencia práctica, si generaliza, es que \textbf{el núcleo nominal es una cota
superior del alcance y hay que medirlo}.

\subsection*{Medición 2 · el corte es exacto, y la capa sí ve lo que la consulta no}

Acá está el resultado central, y la parte que sorprende. Con la misma intervención medimos dos cosas:
cuánto se mueve la \textbf{consulta}, y cuánto se mueve la \textbf{salida de la capa}.

\begin{center}
\begin{tabular}{crr}
\toprule
distancia & consulta & salida de la capa \\
\midrule
$d=1$ & 1.100457 & 0.074570 \\
$d=2$ & 0.579417 & 0.057475 \\
$d=3$ & \textbf{0.000000} & 0.017182 \\
$d=4$ & \textbf{0.000000} & 0.024114 \\
$d=5$ & \textbf{0.000000} & 0.014883 \\
$d=6$ & \textbf{0.000000} & 0.012871 \\
$d=7$ & \textbf{0.000000} & 0.017005 \\
\bottomrule
\end{tabular}
\end{center}

Dos cosas, y la segunda es la que hace que esto no sea obvio.

\textbf{No hay decaimiento.} El movimiento de la consulta pasa de un valor grande a \texttt{0.000000}
de un paso al siguiente, exactamente donde termina el alcance medido.

\textbf{Y el modelo sí ve ese token.} La salida de la capa se mueve a \emph{todas} las distancias,
porque la recurrencia arrastra la información. O sea:

\begin{center}
\emph{El estado ve la secuencia entera; la consulta que lo lee, no.}
\end{center}

Por eso el problema es invisible desde afuera. El modelo \emph{tiene} la información. Lo que no puede
es usarla para decidir qué buscar.

\section{Para qué sirve}

\subsection*{Un diagnóstico que no cuesta nada}

Cambiar un token y ver si la búsqueda se mueve no necesita entrenamiento, ni gradientes, ni etiquetas.
Si no se mueve, ese token es \textbf{invisible} para la búsqueda, y cualquier confianza que el modelo
exprese sobre él \textbf{no está ganada}.

\subsection*{Un arreglo que cuesta el 0,12\,\%}

Ensanchar el núcleo de la convolución que forma la consulta, de 3 a 5, sube el alcance de 2 a 4 y la
parte que faltaba entra en la ventana. Medido en un modelo chico entrenado desde cero, con tres
semillas por condición:

\begin{center}
\begin{tabular}{lcc}
\toprule
 & núcleo 3 & núcleo 5 \\
\midrule
abstención correcta en el caso difícil & $0{,}5850$ -- $0{,}7349$ & $\mathbf{0{,}9931}$ -- $\mathbf{1{,}0000}$ \\
exactitud global & $0{,}91$ -- $0{,}93$ & $\mathbf{0{,}988}$ -- $\mathbf{0{,}993}$ \\
\bottomrule
\end{tabular}
\end{center}

\textbf{Sin un solo solape} entre condiciones, contra un piso trivial de $0{,}4065$. Y cuesta
\textbf{1.024 parámetros sobre 865.651}, o sea el \textbf{0,12\,\%} del modelo: dos taps más, por 128
canales, por 4 bloques.

No es capacidad. Es acceso.

\subsection*{Y no es un caso de borde}

Medido sobre \textbf{33.585 preguntas} de cuatro corpus públicos (SQuAD, Natural Questions, TriviaQA y
HotpotQA), entre el \textbf{90\,\% y el 100\,\%} de las preguntas reales tienen sus partes
discriminantes más separadas que el alcance medido. Y \textbf{empeora con la dificultad}: en las
preguntas multi-hop llega a $1{,}0000$.

La configuración que la tarea sintética fue construida para estudiar es el caso ordinario.

\section{Lo que esto NO afirma}

El límite es parte del resultado, así que conviene decirlo explícito.

\textbf{En un modelo profundo la ventana atenúa, no bloquea.} El corte es exacto en la capa 0, pero
desde la capa 1 la recurrencia restituye la señal atenuada. Medimos la tasa: $1{,}077$ y $1{,}028$ por
token de distancia en modelos de 24 y 48 capas, con $r = 0{,}978$ en los dos. Triplicar los parámetros
y duplicar las capas no cambia la tasa, y eso es lo que hace de la atenuación una propiedad de la
arquitectura y no de un checkpoint.

\textbf{El efecto sobre el comportamiento es transitorio cuando hay capas de sobra.} Al ajustar un
modelo de 24 capas sobre la misma tarea, la diferencia existe y es grande en el paso 100, y
\textbf{desaparece del todo para el paso 400}. Es un resultado negativo, estaba pre-registrado, y lo
informamos.

\begin{center}
\fbox{\begin{minipage}{0.88\textwidth}
\centering
\textbf{Donde no quedan capas con que pagar, la ventana pone un techo.\\ Donde quedan, pone un peaje.}
\end{minipage}}
\end{center}

La afirmación fuerte vale entonces para \textbf{memoria consultada desde una capa temprana}, y no como
predicción de comportamiento para un modelo profundo entrenado hasta converger.

\section{Dónde aplica, según el tamaño}

\begin{center}
\begin{tabular}{lll}
\toprule
 & profundidad & qué pasa \\
\midrule
modelo chico & 4 bloques, lectura en el bloque 0 & \textbf{techo permanente} \\
\texttt{mamba-130m} / \texttt{370m} & 24 y 48 capas & \textbf{peaje}, el efecto se recupera \\
modelos de frontera & decenas de capas, híbridos & \textbf{no medido} \\
\bottomrule
\end{tabular}
\end{center}

Sobre los modelos de frontera lo honesto es esto: \textbf{el mecanismo está presente} —el kernel 4 es
el default en Mamba-2 y en las capas lineales de Qwen3-Next— pero \textbf{que produzca fallas en un
modelo grande entrenado a convergencia no lo medimos y no lo afirmamos}. Dos razones empujan en
contra: si a 24 capas les alcanza para pagar el peaje, es esperable que a 60 les alcance más; y los
modelos de frontera son híbridos, donde una capa de atención completa no tiene esta limitación porque
su consulta ve toda la secuencia.

Donde sí es esperable que importe en un sistema grande es cuando \textbf{la memoria se consulta desde
una capa temprana}, como en una recuperación condicionada por estados de las primeras capas. Ahí no
hay capas aguas arriba para pagar el peaje, y el techo vuelve.

\vspace{1em}
\begin{center}
\fbox{\begin{minipage}{0.9\textwidth}
\textbf{La lección que generaliza.} La consulta con la que se busca en una memoria tiene que formarse
donde ya se vio la pregunta completa. Y se verifica barato en cualquier arquitectura, sin entrenar
nada: cambiá una parte de la consulta y fijate si la búsqueda se mueve. Si no se mueve, esa parte es
invisible.
\end{minipage}}
\end{center}

\end{document}
